\chapter{Conclusioni}
\label{cap:conclusioni}
Da quanto discusso in questo elaborato emerge come il lensing gravitazionale forte rappresenti una delle tecniche osservative più promettenti per lo studio della materia oscura su piccola scala. 
Partendo dal modello cosmologico standard \LCDM, si è visto come la componente dominante della materia dell'Universo sia oscura, ovvero non interagisca con i fotoni, e come differenze nelle sue proprietà si riflettano sulla storia termica dell'Universo e sulla formazione delle strutture. In seguito, si sono riportate le principali evidenze osservative che hanno indotto gli scienziati a teorizzare questa nuova componente cosmica e si sono ripercorse le predizioni teoriche del modello più accreditato, ovvero quello della materia oscura fredda, tra cui il profilo di NFW per la densità di massa degli aloni.
Il modello \LCDM, sebbene rappresenti uno dei più grandi trionfi della cosmologia moderna, presenta delle tensioni su piccola scala rispetto ai risultati osservativi, che potrebbero essere risolte, almeno in parte, da modelli di DM alternativi, come la WDM e la SIDM.
La possibilità di rilevare direttamente o indirettamente sotto-strutture prive di controparte luminosa consente dunque di testare in maniera sempre più stringente le predizioni del modello \LCDM\ e di modelli alternativi.

Tali rilevazioni sono possibili grazie ai fenomeni di lenti gravitazionali.
Sono stati introdotti i rudimenti del lensing gravitazionale, evidenziando il ruolo centrale del potenziale della lente e dei principali parametri geometrici che ne descrivono la fenomenologia.
Le lenti più frequenti in natura sono le ETGs, ovvero galassie ellittiche primordiali, poste a redshift tra 0.1 e 1; per modellarle in maniera semplice si è dunque introdotto il profilo SIE, che permette di riprodurne gli effetti gravitazionali principali, sebbene sia singolare e dunque semplificativo rispetto al profilo reale seguito dalle galassie. 
Attraverso gli anelli di Einstein e gli archi gravitazionali osservabili nelle immagini, è possibile stimare con una precisione dell'1-2\% la massa totale, luminosa e oscura, presente all'interno del raggio di Einstein; tuttavia, tale stima, se effettuata utilizzando il solo lensing, non è univoca. Essa è infatti soggetta a degenerazioni quali la MSD, risolvibili utilizzando, ove possibile, metodi osservativi complementari, come la cinematica stellare.

Combinando quanto presentato nei primi due capitoli, si è mostrato più concretamente come lo strong gravitational lensing possa essere applicato allo studio della materia oscura. Poiché essa non è direttamente osservabile nelle immagini, bisogna sfruttare le sue firme gravitazionali, quali le anomalie di flusso e le perturbazioni alla brillanza superficiale. Nei sistemi reali si formula un problema di inferenza bayesiana gerarchica, che permette, a posteriori, di stimare la probabilità che i parametri del modello potessero generare una specifica immagine. Tuttavia, tale stima è resa difficoltosa dalla forte degenerazione tra la presenza di una sotto-struttura e altre caratteristiche delle immagini riconducibili alle complessità del modello della lente; affinché una rilevazione possa essere considerata robusta, sono necessari più approcci complementari e una particolare attenzione ai modelli utilizzati.

Allo scopo di semplificare, almeno in parte, l'interpretazione di un'osservazione, è utile invertire il problema di inferenza bayesiana e generare simulazioni controllate che mostrino come cambia la morfologia delle immagini di lensing al variare di singoli parametri fisici e geometrici del sistema. Sebbene non consentano di riprodurre tutta la complessità dei sistemi reali, tali simulazioni rappresentano uno strumento utile per costruire un'intuizione fisica del fenomeno e per identificare il legame tra specifiche caratteristiche osservabili e i parametri che le generano. In questo contesto, assumono particolare rilevanza le simulazioni presentate nell'elaborato, progettate per studiare in ambiente controllato l'effetto dei singoli parametri fisici e geometrici sulla morfologia delle immagini. Esse sono state svolte in ambiente ideale, privo di rumore, utilizzando profili semplificati e trascurando sia la risoluzione strumentale che le complessità del modello della lente presentate nei capitoli precedenti.

Future simulazioni più realistiche, che tengano in considerazione quanto qui è stato tralasciato, permetteranno ai futuri fisici e astronomi di comprendere più a fondo e nel dettaglio come la presenza di sotto-strutture oscure possa influenzare la fenomenologia del lensing gravitazionale forte. Ciò si rivelerà uno strumento fondamentale per le presenti e future campagne osservative come \emph{Euclid} e il \emph{Vera Rubin Observatory}, le quali non solo si stima che incrementeranno il numero di lenti finora scoperte, ma grazie alla loro elevata risoluzione angolare permetteranno di abbassare la massa soglia rivelabile. Si stima dunque che la sinergia tra osservazioni e simulazioni permetterà nei prossimi dieci anni di porre vincoli robusti e statisticamente significativi sulla natura della DM.  
